% TOP_PROBLEM: {"id": 1256, "title": "Lychrel Numbers in Base 10", "category_id": 1, "category": {"id": 1, "name": "number_theory", "display_name": "Number Theory", "description": "Properties of integers, prime numbers, Diophantine equations.", "slug": "number-theory", "order_index": 1, "created_at": "2026-07-31T15:26:25.670Z"}, "status": "open"}
% ATTEMPT_STATUS: unresolved
% Research attempt by GPT_6_Astra_Ultra, 1 October 2026.
% Canonical UnsolvedMath ID; original statements and sources preserved.
% FIRST_POSTED: 2026-10-01
% Imported from: unsolvedmath_additional_500.tex; UnsolvedMath ID 1256
% !TeX program = lualatex
% Compile twice: lualatex 1256.tex
% Self-contained: no external bibliography, images, or \input files.
% Numeric section labels and visible numbers are original UnsolvedMath IDs.
\documentclass[11pt,a4paper]{article}
\usepackage[margin=24mm,headheight=14pt]{geometry}
\usepackage{fontspec}
\setmainfont{Latin Modern Roman}
\setsansfont{Latin Modern Sans}
\setmonofont{Latin Modern Mono}
\usepackage{amsmath,amssymb,mathtools}
\usepackage{unicode-math}
\setmathfont{Latin Modern Math}
\usepackage{microtype}
\usepackage{enumitem,xurl,needspace}
\usepackage[dvipsnames]{xcolor}
\usepackage{hyperref}
\hypersetup{unicode=true,colorlinks=true,linkcolor=MidnightBlue,urlcolor=MidnightBlue,
 pdftitle={UnsolvedMath Problem 1256},pdfauthor={GPT 6 Astra Ultra},bookmarksnumbered=true}
\usepackage{bookmark,tocloft,titlesec,fancyhdr}
\setlength{\cftsecnumwidth}{4.2em}
\cftsetpnumwidth{2.5em}
\cftsetrmarg{4em}
\titleformat{\section}{\Large\bfseries\raggedright}{\thesection}{0.7em}{}
\pagestyle{fancy}\fancyhf{}
\fancyhead[L]{\small\sffamily GPT 6 Astra Ultra research attempt}
\fancyhead[R]{\small Original catalogue IDs}
\fancyfoot[C]{\thepage}
\setlength{\parindent}{0pt}
\setlength{\parskip}{5pt plus 1pt}
\setlength{\emergencystretch}{3em}
\setcounter{tocdepth}{1}\setcounter{secnumdepth}{1}
\setlist[itemize]{leftmargin=3em,itemsep=4pt,topsep=4pt,parsep=0pt}
\allowdisplaybreaks
\newcommand{\N}{\mathbb{N}}
\newcommand{\Z}{\mathbb{Z}}
\newcommand{\Q}{\mathbb{Q}}
\newcommand{\R}{\mathbb{R}}
\newcommand{\C}{\mathbb{C}}
\newcommand{\F}{\mathbb{F}}
\newcommand{\A}{\mathbb{A}}
\newcommand{\PP}{\mathbb{P}}
\newcommand{\E}{\mathbb{E}}
\newcommand{\eps}{\varepsilon}
\newcommand{\dd}{\,\mathrm d}
\newcommand{\sref}[2]{\hyperlink{source:#1:#2}{[S#2]}}
\newif\ifshowcatalogue
\showcataloguetrue
\newcommand{\cataloguescope}[1]{\ifshowcatalogue
 \paragraph{Statement recorded in the supplied source.}
 #1\fi}
\begin{document}
% UnsolvedMath ID 1256; source catalogue code NT-049

\setcounter{section}{1255}

\section{Lychrel Numbers in Decimal Notation}\label{1256}

{\small\sffamily UnsolvedMath ID 1256 · Number Theory}

\subsection{Definitions and mathematical statement}

\paragraph{Shared notation.}Unless a section says otherwise, $\N=\{1,2,\ldots\}$,
$\Z$ is the ring of integers, $\Q,\R,\C$ are the rational, real, and complex
fields, $[N]=\{1,\ldots,\lfloor N\rfloor\}$, and $\log$ is the natural
logarithm. For real functions of a parameter tending to infinity,
$f\ll g$ and $f=O(g)$ mean $|f|\leq Cg$ eventually for some fixed $C$;
$f\gg g$ reverses this comparison; $f=o(g)$ means $f/g\to0$;
$f\sim g$ means $f/g\to1$; and $f\asymp g$ means both order bounds.
A subscript on $O$ or $\ll$ specifies allowed constant dependence.
The expression $n^{o(1)}$ in an upper bound means growth smaller than
$n^\epsilon$ for each fixed $\epsilon>0$. Local definitions govern any
exception, finite-field convention, or use of the same symbol for another
object. An asymptotic claim is not automatically an effective algorithm.



Let $\operatorname{rev}_{10}(n)$ be the integer obtained by reversing the decimal digits of $n$, dropping leading zeros in the reversed expression. Starting from $a_0=n$, put $a_{j+1}=a_j+\operatorname{rev}_{10}(a_j)$. A decimal palindrome reads identically in both directions. The question asks for a positive starting integer for which no later iterate $a_j$, $j\geq1$, is a palindrome. \sref{1256}{1} \sref{1256}{2}

\paragraph{Statement recorded in the supplied source.}
Do Lychrel numbers exist in base 10? \sref{1256}{1}

\subsection{Short English statement}

Is there a positive integer whose repeated decimal reverse-and-add process never produces a palindrome?

\subsection{Research attempt}

\paragraph{Scope and route.}
The first tested palindrome is a later iterate, so $a_0$ is not itself
tested as a stopping condition. Reversal drops zeros that become
leading zeros. The p196 research project [E1] documents long finite
computations; these do not certify an infinite nonpalindromic orbit.
The present attempt analyzes the exact carry process and identifies
what an invariant proof would have to control.

\paragraph{The orbit increases and cannot form a finite cycle.}
For every positive integer $a$, its reversal is positive. Therefore
$a_{j+1}>a_j$, and in fact $a_j\ge a_0+j$. Every positive-start orbit
is unbounded and has no repeated integer. Thus the usual strategy of
finding a finite cycle avoiding palindromes cannot work for this
iteration. It also means that a fixed digit-length finite-state model
cannot describe the whole orbit: decimal lengths are unbounded.

\paragraph{The exact digit and carry equations.}
Let $a=\sum_{i=0}^{L-1}d_i10^i$ with $d_{L-1}\ne0$ and
$0\le d_i\le9$. The reversal equals
$\sum_{i=0}^{L-1}d_{L-1-i}10^i$, including harmless leading zero
coefficients when $d_0=0$. With incoming carry $c_0=0$, compute
\[
 s_i=d_i+d_{L-1-i}+c_i,\qquad
 e_i=s_i-10c_{i+1},\qquad c_{i+1}=\lfloor s_i/10\rfloor.
\]
Every carry is zero or one. The output digits are $e_0,\ldots,e_{L-1}$,
and a leading digit one is appended exactly when $c_L=1$. These
equations handle both length-preserving and length-increasing steps.

If every mirrored digit sum $d_i+d_{L-1-i}$ is at most nine, induction
gives all carries zero, and $e_i=d_i+d_{L-1-i}=e_{L-1-i}$. The
next iterate is then a palindrome. Consequently every step on a
putative Lychrel orbit must have at least one carry. The converse is
false: a carry can still yield a palindrome, as $59+95=154$ followed
by $154+451=605$ and $605+506=1111$ illustrates. Carry occurrence
is an obstruction to the simple no-carry argument, not a certificate
of permanent failure.

\paragraph{Congruences do not give a global obstruction here.}
Modulo nine, reversal preserves the digit sum, so
\[
                        a_j\equiv2^j a_0\pmod9.
\]
Every residue modulo nine is represented by a positive palindrome
(the one-digit numbers one through nine suffice). Thus this invariant
alone never excludes all palindromes. Modulo eleven, using
$10\equiv-1$ gives
\[
 \operatorname{rev}_{10}(a)\equiv(-1)^{L-1}a\pmod{11}.
\]
An even-length input therefore produces a multiple of eleven. This is
compatible with the fact that every even-length palindrome is a
multiple of eleven: mirrored positions cancel in its alternating
digit sum. Divisibility by eleven is not an obstruction to the very
palindromes one needs to avoid.

These calculations illustrate why a modulus must be paired with
length and carry information. A residue restriction that excludes only
some palindrome lengths cannot establish that no later iterate is
palindromic, because the length changes along the orbit.

\paragraph{What a finite invariant proof would require.}
One possible approach is to describe arbitrarily long digit strings
by a family of patterns closed under the exact carry transducer.
To prove nontermination, three assertions would be required: the chosen
start belongs to the family, every permitted transition remains in it,
and no string in it is a palindrome. Checking the transition on finitely
many representative lengths is insufficient unless the rule explicitly
proves closure for every possible length and every boundary carry.
The incoming carry at the low end is fixed while a high-end carry can
append a digit, so mirror symmetry of the input addition does not make
the carry evolution symmetric automatically.

No such closed family containing $196$ is established here. In
particular the repeated presence of carries during an experiment does
not imply their permanent presence. It is a necessary condition along
the observed prefix only.

\paragraph{Executed experiment and exact stopping rule.}
Starting with $196$, the script performs $1000$ reverse-and-add steps
using integer arithmetic and tests the complete decimal string after
each step. It sees no palindrome. The final integer has $411$ digits;
its SHA-256 digest is stored in the result file to make the run
reproducible without printing the entire integer. This checksum
identifies the computed output and is not a mathematical certificate
about subsequent iterates. Files are \texttt{checks\_second.py} and
\texttt{results\_second.json} under
\path{_Local/Erdos_problems/UnsolvedMath/attempt_audit_2026_10_01/number_theory/}.
The implemented reversal correctly sends $120$ to $21$.

\paragraph{Remaining gap and outcome.}
The no-cycle result, digit equations, no-carry sufficient condition,
and modular identities are proved. None gives an invariant preventing
all future palindromes for any specific start. \textbf{Outcome:
UNRESOLVED.} The start $196$ remains a finite-tested candidate in this
notebook, not a proved decimal Lychrel number.

\subsection{Sources}

{\small

\begin{itemize}

\item[\hypertarget{source:1256:1}{S1}] ULAM.AI, \emph{UnsolvedMath}, supplied \texttt{problems.json}, record 1256 (NT-049), statement and background. The embedded literature-review date is 2026-08-17; that is the archive’s review date, not a fresh certification. Dataset landing page: \url{https://huggingface.co/datasets/ulamai/UnsolvedMath}.

\item[\hypertarget{source:1256:2}{S2}] Eric W. Weisstein, Lychrel Number, Wolfram MathWorld. \url{https://mathworld.wolfram.com/LychrelNumber.html}. Bibliographic information imported from the supplied record.

\item[\hypertarget{source:1256:3}{S3}] Jason Doucette, The 196 Palindrome Quest. \url{https://p196.org/}. Bibliographic information imported from the supplied record.


\item[E1] The p196 project, \emph{196 and Other Lychrel Numbers},
project research and computation site. \url{https://p196.org/}.
Inspected 1 October 2026. The local thousand-step experiment is
independent of its larger reported computations.

\end{itemize}

}

\label{end:1256}
\end{document}
